\documentclass[11pt]{article}
\usepackage[margin=1in]{geometry}
\usepackage{amsmath,amssymb,amsthm}
\usepackage{xurl}
\let\oldus\_
\renewcommand{\_}{\oldus\allowbreak}
\newtheorem{lemma}{Lemma}
\newcommand{\Z}{\mathbb{Z}}
\newcommand{\Zp}{\mathbb{Z}_p}

\title{A disproof of Conjecture 00000000781\\[2pt]
\large The closed-subgroup lattice of $1+p\Zp$}
\date{}

\begin{document}
\sloppy
\maketitle

\section{The conjecture and the reading}

\begin{quote}
\emph{Definition:} The lattice of closed subgroups of $1+p\Zp$ is the combinatorial structure of the
multiplicative group. \emph{Conjecture:} This subgroup lattice is isomorphic to a quotient lattice of the
integers; the isomorphism is the valuation stratification; and the isomorphism preserves the stratified
structure of the valuation filtration.
\end{quote}

\textbf{Reading.} Neither the English nor the Chinese text restricts $p$ (in particular, neither says
``odd''), so the statement is read for every prime $p$; we refute it at $p=2$. The conclusion is a
conjunction, so it suffices to refute the first clause, the existence of an isomorphism. The lattice is the set
of closed subgroups of $U_1(p)=\{u\in\Zp^\times : p \mid u-1\}$ ordered by inclusion (the topology is the
subspace topology from the unit group); a lattice isomorphism is the same as an order isomorphism.

\textbf{Convention for ``quotient lattice of the integers''.} This is not a defined term. We formalize four
readings (Lean: \texttt{C781.QuotientLatticeOfInt}): a lattice $L$ that is a surjective lattice-homomorphic
image of (a) the chain $(\Z,\le)$, (b) the subgroup lattice of $\Z$, (c) its order dual (the divisibility
lattice), or (d) a lattice isomorphic to the subgroup lattice of $\Z/n\Z$. Other senses are not covered,
and the statement for odd $p$ (where it is plausibly true) is not addressed.

\section{Proof}

\begin{lemma}[Distributivity]
Every lattice in (a)--(d) is distributive.
\end{lemma}
\begin{proof}
A surjective lattice homomorphic image of a distributive lattice is distributive. The chain is distributive.
Every subgroup of $\Z$ is $n\Z$, with $n\Z\cap m\Z=\mathrm{lcm}(n,m)\Z$ and $n\Z+m\Z=\gcd(n,m)\Z$; the
identity $\gcd(\mathrm{lcm}(a,b),\mathrm{lcm}(a,c))=\mathrm{lcm}(a,\gcd(b,c))$ is checked on prime
factorizations, giving distributivity of the subgroup lattice of $\Z$, hence of its dual. For
$\Z/n\Z$, pulling back along the surjection $\Z\to\Z/n\Z$ is an order embedding preserving meets and joins into
the subgroup lattice of $\Z$. (Lean: \texttt{distrib\_of\_surj}, \texttt{dual\_distrib}, \texttt{gcd\_lcm\_key},
\texttt{subZ\_distrib}, \texttt{subZMod\_distrib}, \texttt{quotient\_distrib}.)
\end{proof}

\begin{lemma}[A diamond at $p=2$]
Let $\rho:\Z_2\to\Z/8$ be reduction and, for $a\in\{3,5,7\}$, $H_a=\{u\in U_1(2):\rho(u)\in\{1,a\}\}$. Then
$H_a$ is a closed subgroup, $H_7\cap H_3=H_7\cap H_5=\{\rho=1\}$, and the only closed subgroup containing
$H_7$ and $H_a$ ($a=3,5$) is $U_1(2)$.
\end{lemma}
\begin{proof}
$\rho$ is continuous (its fibres are balls of radius $2^{-3}$) into a discrete space, so $H_a=\rho^{-1}\{1,a\}$
is closed; it is a subgroup because $a^2=1$ in $\Z/8$. Every $u\in U_1(2)$ has $\rho(u)=1+2s$, so
$\rho(u)\in\{1,3,5,7\}$. If a closed subgroup $X$ contains $H_7$ and $H_a$, then it contains $-1$
($\rho(-1)=7$); if $\rho(u)$ is neither $1,7$ nor $a$, then $\rho(-u)=a$, so $-u\in X$ and $u=(-1)(-u)\in X$.
(Lean: \texttt{pi8\_cont}, \texttt{Hres}, \texttt{top\_of}, \texttt{isLUB\_pair}, \texttt{isGLB\_pair}.)
\end{proof}

\begin{lemma}[Main theorem]
There is no order isomorphism from the closed subgroups of $1+2\Z_2$ onto a lattice as in (a)--(d).
\end{lemma}
\begin{proof}
Suppose $e$ is one. An order isomorphism carries meets and joins (as greatest lower bounds and least upper
bounds) to meets and joins, so by Lemma 2, $e(H_7)\wedge e(H_3)=e(H_7)\wedge e(H_5)$ and
$e(H_7)\vee e(H_3)=e(H_7)\vee e(H_5)$. Since $L$ is distributive (Lemma 1), cancellation gives
$e(H_3)=e(H_5)$, so $H_3=H_5$. But $3$ is a unit of $\Z_2$ with $3-1=2$, so $3\in H_3\setminus H_5$. Contradiction.
\end{proof}

\section{Lean and audit}

File \texttt{Conjecture781/Basic.lean} (identical to \texttt{Tlmc/C781.lean}). The statement block defines
\texttt{U1}, \texttt{QuotientLatticeOfInt}, \texttt{Conjecture p} (there is a lattice $L$ with
\texttt{QuotientLatticeOfInt L} and an order isomorphism \texttt{ClosedSubgroup (U1 p)} $\simeq_o L$) and
\texttt{Claim := not (forall p prime, Conjecture p)}. The theorem is \texttt{C781.main : C781.Claim}.
Axioms: only \texttt{propext}, \texttt{Classical.choice}, \texttt{Quot.sound}; no \texttt{sorry}.

\textbf{Python check.} The same diamond exists in $(\Z/8)^\times$: subgroups $\{1,3\},\{1,5\},\{1,7\}$ meet
pairwise in $\{1\}$ and join to everything (attempt file \texttt{check\_code}).

\end{document}
